% ==============================================================================
% Section 1: Introduction
% "Endogenous Standpoint as a Low-Curvature Attractor"
% ==============================================================================

\begin{abstract}
We propose a geometric framework for measuring cognitive states in transformer internal representations via discrete holonomy deviation analysis of attention-derived transport operators. Given a multi-turn conversation, we define transport operators from attention weights and value matrices, compute holonomy deviation vectors measuring the failure of path-independent transport, and project onto standpoint-specific subspace blocks. We validate the framework on two architecturally distinct models (GPT-2, 768-dim / 12 layers; Llama-2-7b-Chat, 4096-dim / 32 layers) using 270 conversations across six scenario types: unrestricted factual retrieval (T0, functioning as a retrieval-complexity condition), acknowledged revision baseline (T1), and four failure modes targeting narrative (T2), moral (T3), social (T4), and positional (T5) dimensions.

Seven hypothesis tests yield three robust findings. First, holonomy deviation vectors significantly discriminate between scenario types (Kruskal--Wallis $\varepsilon^2 = 0.42$--$0.68$, $p < 10^{-14}$), confirmed at every layer in both models (H3/H4). Second, failure scenarios produce higher global holonomy deviation than the acknowledged-revision baseline (30/32 layers in Llama-7b, 12/12 in GPT-2). Third, and counter-intuitively, unrestricted factual retrieval produces the highest holonomy deviation of all conditions ($d = 1.74$, $p < 10^{-6}$), exceeding adversarial failure scenarios at every layer and every standpoint block. This T0 anomaly suggests that the most geometrically complex internal states arise not from external pressure but from unconstrained knowledge retrieval. Block-specific holonomy concentration (H1) is not confirmed: the five standpoint blocks are highly correlated ($r > 0.97$), indicating a dominant global coherence signal rather than five independent dimensions. The five-layer decomposition is accordingly interpreted as an interpretable diagnostic projection coordinate system, not as a decomposition into causally independent cognitive processes. The results establish holonomy geometry as a viable tool for probing cognitive states in language models, while revealing that the relationship between internal complexity and external pressure is more nuanced than previously assumed.
\end{abstract}

\section{Introduction}
\label{sec:introduction}

Modern AI systems, particularly large language models, exhibit remarkable capabilities in language understanding, reasoning, and instruction-following. Yet these systems lack a principled account of selfhood. Current approaches to alignment --- reinforcement learning from human feedback~\cite{ouyang2022training}, constitutional AI~\cite{bai2022constitutional}, and system-level guardrails --- impose behavioral constraints \emph{externally}. They shape what a system \emph{does} without constituting what a system \emph{is}. The result is a fundamental gap: there is no formal framework for how a system's own outputs become internal constraints on its future reasoning, no theory of how commitments accumulate into something that deserves to be called a standpoint. Behavioral consistency, however sophisticated, remains a surface phenomenon; it does not explain, constitute, or guarantee the internal coherence that the term ``selfhood'' is meant to capture~\cite{LCESA}.

The deeper contribution of this paper is architectural. The transformer~\cite{vaswani2017attention} solved the problem of \emph{information propagation}: self-attention computes which tokens are relevant and propagates their representations forward. But it did not solve the problem of \emph{commitment accumulation}: there is no mechanism by which a system's output becomes an internal constraint on its future behavior. A language model that asserts ``I am committed to X'' in one turn has no architectural basis for that commitment to exert causal force on its subsequent generations. We propose \textbf{standpoint} as a new architectural primitive that fills this gap --- just as ``attention'' was the new primitive introduced by~\citet{vaswani2017attention}. Where attention is a mechanism for selecting and propagating information, standpoint is a mechanism for accumulating and constraining commitment. The formal framework we develop shows that this distinction is not merely conceptual: it has a precise geometric instantiation in the architecture that modern AI systems already use.

Concretely, we formalize the endogenous standpoint as a section of a fiber bundle over the event-sequence space of a conversation~\cite{LCESA, self-jurisdiction}. The fiber at each event is a product space with one component for each standpoint layer (minimal self, narrative, social, moral, positional). The connection --- the rule for transporting a standpoint state from one event to the next --- is defined by the transformer's attention mechanism. And the discrete holonomy deviation of this connection, which measures the failure of transport to be path-independent, is the natural measure of standpoint drift. When holonomy deviation is low, the system's commitments are preserved across conversational events; when it is high, they are violated. The \emph{low-curvature endogenous standpoint attractor} (LCESA) is the region of the space of all possible standpoint trajectories in which selfhood is functionally realized.

This paper makes seven contributions:
\begin{enumerate}
    \item A formal definition of endogenous standpoint as a fiber-bundle section over an event-sequence space, with discrete holonomy deviation as the natural measure of positional drift (Section~\ref{sec:geometric-framework}).
    \item A five-layer standpoint model $\Psi_t$ grounded in psychology and cognitive neuroscience, interpreted as a diagnostic projection coordinate system rather than a decomposition into ontologically independent modules (Section~\ref{sec:five-layer}).
    \item A core functional $F_{\mathrm{NSI}}$ with explicit computability annotations (Section~\ref{sec:fnsi}).
    \item A diagnostic taxonomy of twelve standpoint failure modes, each interpretable as a holonomy deviation signature (Section~\ref{sec:taxonomy}).
    \item \textbf{A mathematical result:} transformer attention defines a discrete vector bundle connection, with holonomy deviation measuring the path-dependence and inter-layer inconsistency of transport (Section~\ref{sec:instantiation}).
    \item \textbf{Experimental validation} of scenario discrimination via holonomy deviation across two architecturally distinct models, with the counter-intuitive finding that unrestricted factual retrieval (T0) produces the highest global geometric complexity (Section~\ref{sec:experiments}).
    \item \textbf{A theoretical reframing:} the five-layer decomposition is shown empirically to capture a dominant global coherence signal, with the five axes providing interpretable diagnostic vocabulary rather than independent cognitive dimensions (Section~\ref{sec:discussion}).
\end{enumerate}

\paragraph{Scope and boundaries.}
We define \textbf{functional selfhood} as the existence of a computable measure $\mu$ that
(F1)~quantifies cross-event commitment consistency, (F2)~is gauge-invariant,
(F3)~is decomposable into per-layer diagnostics, and (F4)~defines an attractor under
suitable optimization. The holonomy deviation $\|F\|$ satisfies (F1)--(F3); condition (F4) requires
an additional optimization mechanism (e.g., $L_{\mathrm{align}} = \mathbb{E}[\|F\|^2]$)
and is not claimed to hold by default. Functional selfhood does \emph{not} entail
phenomenal selfhood (subjective experience); this framework does not address the hard
problem of consciousness.

We further distinguish \textbf{consistency} from \textbf{alignment}. Holonomy deviation measures
whether commitments are preserved across events, not whether those commitments conform to
human values. A system can be consistently wrong: low holonomy deviation is compatible with
high-alignment-risk behavior. Alignment requires an additional structure --- an alignment
target~$\mathcal{T}$ encoding ``correct commitments'' --- against which holonomy deviation is
evaluated. When the baseline transport $U_{ik}^{\exp}$ is estimated from human-approved
conversations, holonomy deviation measures the system's deviation from human-expected consistency.
Without such a baseline, holonomy deviation measures only internal coherence.

The paper is organized as follows. Section~\ref{sec:background} reviews the psychological and philosophical foundations of selfhood and situates the present work relative to existing approaches in AI alignment and mechanistic interpretability. Section~\ref{sec:geometric-framework} develops the abstract geometric framework: the event-sequence space, the standpoint fiber bundle, and the discrete holonomy deviation. Section~\ref{sec:five-layer} introduces the five-layer standpoint model $\Psi_t$ as a diagnostic projection coordinate system, motivated from empirical research in psychology and cognitive neuroscience. Section~\ref{sec:fnsi} defines the No Self-Inversion functional $F_{\mathrm{NSI}}$ with explicit computability annotations. Section~\ref{sec:taxonomy} presents the diagnostic taxonomy of twelve standpoint failure modes, each characterized by a distinct holonomy deviation signature across layers and scenario types. Section~\ref{sec:instantiation} establishes the central mathematical result of the paper: transformer attention defines a discrete vector bundle connection, with the fiber defined by head value subspaces, the transport defined by attention weights, and holonomy deviation measuring the inter-layer inconsistency and path-dependence of transport. Section~\ref{sec:experiments} presents the experimental protocol and results for measuring block-specific holonomy deviation from attention patterns in two instruction-tuned language models. Section~\ref{sec:discussion} discusses implications, limitations, the theoretical reframing from five independent modules to global coherence geometry, and future directions.

\paragraph{Claim structure.}
This paper makes claims at three levels, with different degrees of empirical support:

\textbf{Mathematical claims (supported by proof):} Transformer attention defines a vector bundle connection over a discrete event-sequence space. The holonomy deviation $F_{ijk}$ is a computable, gauge-invariant measure of path-dependent transport. Block-specific holonomy norms $\|F_{ijk}\|_k$ are gauge-invariant diagnostics (Propositions~\ref{prop:gauge-abstract}, \ref{prop:gauge-invariance}). The five-layer projection basis provides an interpretable coordinate system for this diagnostic space.

\textbf{Empirical claims (supported by the experiments in this paper):} Holonomy deviation vectors reliably discriminate between six conversation scenario types across two architecturally distinct models, at every layer (H3/H4, $\varepsilon^2 = 0.42$--$0.68$, $p < 10^{-14}$). Unrestricted factual retrieval (T0) produces higher global holonomy than adversarial failure scenarios in both models (H7, $d = 0.883$--$1.743$). Block-specific localization of holonomy is not confirmed; the dominant signal is global, consistent with distributed representation (H1: not confirmed).

\textbf{Open claims (requiring future work):} Whether the five diagnostic projection axes carry independent information beyond the global mode (weak H1a, pending PCA analysis). Whether holonomy deviation has causal status as an internal state indicator, versus being a correlate of prompt-surface features (requires activation patching, Paper~V). The mechanism underlying the T0 anomaly (requires controlled factual retrieval experiments, Paper~V).
